\subsection{An explicit cubic inertia witness}

The following construction gives a local arithmetic witness for the
strictness of the second odd CM restriction.

\begin{proposition}\label{prop:five-inertia-odd-cubic}
There is an element \(w\in U_D\cap H\) whose full harmonic
frame is one and whose cubic frame at five is
\[
 T_5(w)=\left(\frac{1+5\omega}{1+5\omega^2}\right)^{-2}.
\]
This value has infinite pro-five order.
\end{proposition}
\begin{proof}
Use the CM curve \(E_3:y^2=x^3+1\).
It has good reduction at five.  The affine point counts
above \(x=0,\ldots,4\) are \(2,0,2,0,1\), so
\(\#E_3(\mathbf F_5)=6\).  The five-power Frobenius
therefore satisfies \(\pi^2=[-5]\) on the reduction.
Over the unramified quadratic extension
\(K=\mathbf Q_5(\omega)\), with integers
\(O=\mathbf Z_5[\omega]\), the formal CM parameter
\(t=-x/y\) has \(\omega\)-action \(t\mapsto\omega t\).
As in the preceding local computation,
\[
 [-5](t)\equiv-5t\pmod{t^2},\qquad
 [-5](t)\equiv t^{25}\pmod5.
\]
The formal group is the Lubin--Tate module for the
uniformizer \(-5\), and all five-power torsion lies in
it.  Hence inertia \(I_5\) has image \(O^\times\) on
its rank-one \(O\)-Tate module
\cite[Theorem~2(b),(d),(e), pp.~383--384]{LubinTate1965}.

Choose
\[
                       a=\frac{1+5\omega}{1+5\omega^2}.
\]
It has norm one and belongs to \(1+5O\).  Its logarithm
has leading term \(5(\omega-\omega^2)\), of valuation
one, so it has infinite order.  Choose \(j\in I_5\)
whose cubic Tate action is multiplication by \(a\).
The determinant and the Weil pairing give
\(\lambda_5(j)=1\), and prime-to-five roots of unity
are unramified.  Thus \(\lambda(j)=1\) at every prime.

Choose \(h\in I_3\) with \(\lambda_3(h)=-1\);
all its other cyclotomic components are one.  Such an
element exists in the local cyclotomic inertia image.
It sends \(\omega\) to \(\omega^2\).  Accordingly its
cubic Tate action conjugates multiplication by \(a\)
to multiplication by \(\bar a\).
Put \(z=[h,j]=h j h^{-1}j^{-1}\).
The tangential pro-\(p\) action itself is unramified
outside \(p\), as in \cite[\S8.5]{HainMatsumoto2003}.
For every \(p\), one of the two factors has trivial
pro-\(p\) action.  Therefore
\[
                    \lambda(z)=1,\qquad \Phi_p(z)=1
                                      \quad\hbox{for every }p.
\]
The rational Kummer cocycle formula gives
\[
                \rho_q(z)=(\lambda(h)-1)\rho_q(j)
                                      \quad\hbox{for every }q.
\]
Each such value is supported only at the coefficient
prime three.  Let \(e_5\in\widehat{\mathbf Z}\)
be the idempotent supported at five, and put
\(w=z^{e_5}\), using the continuous power map on
\(\overline{\langle z\rangle}\).
Additivity on multiplier one gives \(\rho_q(w)=0\)
for every rational prime \(q\).  Thus
\(w\in P_D\cap H\), and arithmetic vanishing of the
cubic residue defect gives \(w\in T_D\cap H\).

We check that the harmonic image is already exactly
trivial.  Both inertia groups \(I_3\) and \(I_5\)
fix \(\mathbf Q(i)\), which is unramified at these
primes.  The Gaussian Tate representation restricted
to \(G_{\mathbf Q(i)}\) is abelian: all CM
endomorphisms are defined there, so the image lies in
the commutative unit group
\((\mathbf Z[i]\otimes\widehat{\mathbf Z})^\times\).
Consequently \(\rho_{E_4}(z)=1\), and likewise
\(\rho_{E_4}(w)=1\).  The signed harmonic identification
on the cyclotomic arithmetic subgroup yields
\[
                               W(w)=1.
\]
In particular \(w\in U_D\cap H\).

On the cubic Tate module at five,
\(\rho_{E_3,5}(z)=\bar a/a=a^{-2}\).
This value is pro-five, so the \(e_5\) power preserves it.
The signed cubic frame on arithmetic multiplier one
is \(T=\kappa\rho_{E_3}\), where
\(\kappa\) is the fourth-root Kummer character of
three.  Since \(w\in D\), it fixes that fourth root,
and \(\kappa(w)=1\).  It follows that
\[
                              T_5(w)=a^{-2}.
\]
The logarithm calculation above proves its infinite pro-five order.
\end{proof}
